% part: intuitionistic-logic
% chapter: introduction
% section: axiomatic-derivations

\documentclass[../../../include/open-logic-section]{subfiles}

\begin{document}

\olfileid[hi]{int}{int}{axd}

\olsection{अभिगृहीतीय \usetoken{P}{derivation}}

अंतःप्रज्ञावादी प्रतिज्ञप्तिक तर्कशास्त्र की अभिगृहीतीय !!{derivation}याँ
संकल्पनात्मक रूप से सबसे सरल और ऐतिहासिक रूप से सबसे पहली
!!{derivation}-प्रणालियाँ हैं। वे ठीक शास्त्रीय प्रतिज्ञप्तिक तर्कशास्त्र की
तरह कार्य करती हैं।

\begin{defn}[!!{derivability}]
यदि $\Gamma$, $\Lang L$ के !!{formula}ों का समुच्चय है, तो $\Gamma$ से कोई
\emph{!!{derivation}}, !!{formula}ों का ऐसा परिमित अनुक्रम $!A_1$,
\dots,~$!A_n$ है जिसमें प्रत्येक $i \le n$ के लिए निम्न में से कोई एक
शर्त पूरी होती है:
\begin{enumerate}
\item $!A_i \in \Gamma$; अथवा
\item $!A_i$ कोई अभिगृहीत है; अथवा
\item $!A_i$, किन्हीं $!A_j$ और $!A_k$ से, जहाँ $j < i$ और $k < i$,
  मोडस पोनेन्स द्वारा मिलता है; अर्थात $!A_k \ident !A_j \lif !A_i$।
\end{enumerate}
\end{defn}

\begin{defn}[अभिगृहीत]
अंतःप्रज्ञावादी प्रतिज्ञप्तिक तर्कशास्त्र के \emph{अभिगृहीतों} का समुच्चय
$\PAx$, निम्न रूपों वाले सभी !!{formula}ों से बनता है:
\begin{align}
  & (!A \land !B) \lif !A \ollabel{ax:land1}\\
  & (!A \land !B) \lif !B \ollabel{ax:land2}\\
  & {}!A \lif (!B \lif (!A \land !B)) \ollabel{ax:land3}\\
  & {}!A \lif (!A \lor !B) \ollabel{ax:lor1}\\
  & {}!A \lif (!B \lor !A) \ollabel{ax:lor2}\\
  & (!A \lif !C) \lif ((!B \lif !C) \lif ((!A \lor !B) \lif !C)) \ollabel{ax:lor3}\\
  & {}!A \lif (!B \lif !A) \ollabel{ax:lif1}\\
  & (!A \lif (!B \lif !C)) \lif ((!A \lif !B) \lif (!A \lif !C)) \ollabel{ax:lif2}\\
  & \lfalse \lif !A \ollabel{ax:lfalse1}
\end{align}
\end{defn}

\begin{defn}[!!{derivability}]
कोई !!{formula}~$!A$, $\Gamma$ से \emph{!!{derivable}} है, जिसे
$\Gamma \Proves !A$ लिखते हैं, यदि $\Gamma$ से ऐसी कोई !!{derivation} है
जो $!A$ पर समाप्त होती है।
\end{defn}

\begin{defn}[प्रमेय]
कोई !!{formula}~$!A$ \emph{प्रमेय} है, यदि रिक्त समुच्चय से~$!A$ की कोई
!!{derivation} है। यदि $!A$ प्रमेय है तो हम $\Proves !A$, और यदि नहीं है
तो $\Proves/ !A$ लिखते हैं।
\end{defn}

\begin{prop}
  यदि अंतःप्रज्ञावादी तर्कशास्त्र की अभिगृहीतीय !!{derivation}-प्रणाली में
  $\Gamma \Proves !A$, तो शास्त्रीय तर्कशास्त्र में भी
  $\Gamma \Proves !A$। विशेषतः, यदि $!A$ अंतःप्रज्ञावादी प्रमेय है, तो वह
  शास्त्रीय प्रमेय भी है।
\end{prop}

\begin{proof}
  प्रत्येक अंतःप्रज्ञावादी अभिगृहीत शास्त्रीय अभिगृहीत भी है; अतः
  अंतःप्रज्ञावादी तर्कशास्त्र की प्रत्येक !!{derivation}, शास्त्रीय
  तर्कशास्त्र की भी !!{derivation} है।
\end{proof}

\end{document}
